\documentclass[12pt]{article}
\usepackage[utf8]{inputenc}
\usepackage[T1]{fontenc}
\usepackage[margin=1in]{geometry}
\usepackage{setspace}
\usepackage{xcolor}
\usepackage[colorlinks=true, urlcolor=blue!60!black, linkcolor=blue!60!black]{hyperref}
\setlength{\parindent}{2em}
\setlength{\parskip}{1.25em}
\renewcommand{\baselinestretch}{1.0}
\begin{document}
\begin{center}
{\LARGE \textit{Gibbons v. Ogden}, 22 U.S. (9 Wheat.) 1 (1824)\par}
\vspace{0.8em}
{\large Case Brief\par}
\end{center}
\vspace{1em}

\section{Memory Jogger}
Congress's Commerce Clause power encompasses interstate navigation and is plenary within its sphere, so a federal coasting license preempts New York's conflicting steamboat monopoly under the Supremacy Clause.

\section{Detailed Case Facts}
Beginning in 1787 (with an initial grant to John Fitch) and continuing through statutes enacted in 1798, 1803, 1808, and 1811, the New York Legislature granted Robert R. Livingston and Robert Fulton the exclusive right to navigate all waters within New York's territory and jurisdiction using vessels propelled by fire or steam for a term extending up to thirty years, enforced by summary forfeiture of unlicensed steamboats (\href{https://biotech.law.lsu.edu/cases/pp/Gibbons.htm}{LSU Law Center text of \textit{Gibbons v. Ogden}}; \href{https://www.oyez.org/cases/1789-1850/22us1}{Oyez}). Livingston and Fulton assigned the exclusive right to operate steamboats between Elizabethtown (and other points in New Jersey) and New York City to John R. Livingston, who in turn assigned that franchise to Aaron Ogden, a former Governor of New Jersey (\href{https://biotech.law.lsu.edu/cases/pp/Gibbons.htm}{LSU Law Center}; \href{https://history.nycourts.gov/wp-content/uploads/2026/05/25_Judicial_Notice-compressed.pdf}{Historical Society of the New York Courts}).

Thomas Gibbons, a New Jersey citizen and former business partner of Ogden, broke with Ogden and began operating two rival steamboats---the \textit{Stoudinger} and the \textit{Bellona}---carrying passengers between Elizabethtown, New Jersey, and New York City (\href{https://biotech.law.lsu.edu/cases/pp/Gibbons.htm}{LSU Law Center}; \href{https://www.oyez.org/cases/1789-1850/22us1}{Oyez}). Gibbons held no license from Livingston, Fulton, or Ogden, but his vessels were enrolled and licensed to engage in the coasting trade under the federal Coasting Act of February 18, 1793 (``An act for enrolling and licensing ships and vessels to be employed in the coasting trade and fisheries, and for regulating the same'') (\href{https://biotech.law.lsu.edu/cases/pp/Gibbons.htm}{LSU Law Center}). Meanwhile, neighboring states had enacted retaliatory measures---Connecticut barred vessels licensed under New York's monopoly from entering Connecticut waters, and New Jersey passed an ``act of retortion'' authorizing damages and treble costs against anyone enforcing New York's monopoly against New Jersey citizens---creating an escalating interstate commercial conflict (\href{https://opencasebook.org/casebooks/363-constitutional-law/resources/2.3.1-gibbons-v-ogden/}{OpenCasebook}). In October 1818, Ogden filed a bill in equity in New York's Court of Chancery seeking an injunction to stop Gibbons from operating his steamboats in New York waters (\href{https://opencasebook.org/casebooks/363-constitutional-law/resources/2.3.1-gibbons-v-ogden/}{OpenCasebook}).

\section{Procedural History}
\begin{enumerate}
\item \textbf{New York Court of Chancery (\textit{Ogden v. Gibbons}, 4 Johns. Ch. 150 (N.Y. Ch. 1819)):} Chancellor James Kent granted and later made perpetual an injunction restraining Gibbons from navigating New York waters by steam, holding that New York's monopoly statutes were valid exercises of concurrent state authority and that the federal Coasting Act license merely conferred American national character and customs privileges rather than an affirmative right to navigate state waters (\href{https://history.nycourts.gov/wp-content/uploads/2026/05/25_Judicial_Notice-compressed.pdf}{Historical Society of the New York Courts}; \href{https://biotech.law.lsu.edu/cases/pp/Gibbons.htm}{LSU Law Center}).
\item \textbf{New York Court for the Trial of Impeachments and Correction of Errors (\textit{Gibbons v. Ogden}, 17 Johns. 488 (N.Y. 1820)):} New York's highest court of last resort unanimously affirmed Chancellor Kent's decree (\href{https://history.nycourts.gov/wp-content/uploads/2026/05/25_Judicial_Notice-compressed.pdf}{Historical Society of the New York Courts}).
\item \textbf{U.S. Supreme Court---First Appeal (\textit{Gibbons v. Ogden}, 19 U.S. (6 Wheat.) 448 (1821)):} Gibbons appealed from the interlocutory refusal to dissolve the injunction, but the Supreme Court dismissed the appeal for want of jurisdiction because the record did not show a final decree as required by \S~25 of the Judiciary Act of 1789 (\href{https://supreme.justia.com/cases/federal/us/19/448/}{Justia}).
\item \textbf{U.S. Supreme Court---Final Appeal (\textit{Gibbons v. Ogden}, 22 U.S. (9 Wheat.) 1 (1824)):} Following entry of a final decree, Gibbons appealed again. Argued February 5--9, 1824 (with Daniel Webster and U.S. Attorney General William Wirt representing Gibbons, and Thomas J. Oakley and Thomas Addis Emmet representing Ogden) and decided March 2, 1824, the Supreme Court reversed and annulled the New York decree and ordered Ogden's bill dismissed (\href{https://www.oyez.org/cases/1789-1850/22us1}{Oyez}; \href{https://ballotpedia.org/Gibbons_v._Ogden}{Ballotpedia}; \href{https://www.courtlistener.com/opinion/85412/gibbons-v-ogden/}{CourtListener}).
\end{enumerate}

\section{Judicial Votes}
\textbf{Vote Breakdown:} 6--0 unanimous decision of the participating Justices (one Justice not participating) (\href{https://ballotpedia.org/Gibbons_v._Ogden}{Ballotpedia}; \href{https://www.oyez.org/cases/1789-1850/22us1}{Oyez}).
\begin{itemize}
\item \textbf{Majority (5):} Chief Justice John Marshall (author), joined by Associate Justices Bushrod Washington, Thomas Todd, Gabriel Duvall, and Joseph Story (\href{https://teachingamericanhistory.org/document/gibbons-v-ogden/}{Teaching American History}; \href{https://ballotpedia.org/Gibbons_v._Ogden}{Ballotpedia}).
\item \textbf{Concurring in the Judgment (1):} Associate Justice William Johnson (separate opinion) (\href{https://teachingamericanhistory.org/document/gibbons-v-ogden/}{Teaching American History}).
\item \textbf{Dissenting (0):} None.
\item \textbf{Not Participating (1):} Associate Justice Smith Thompson did not participate in the discussion or decision of the case (\href{https://www.oyez.org/cases/1789-1850/22us1}{Oyez}).
\end{itemize}

\section{Holding}
\textbf{Judgment:} Reversed and annulled the decree of the New York Court for the Trial of Impeachments and Correction of Errors; directed that Aaron Ogden's bill in equity be dismissed (\href{https://biotech.law.lsu.edu/cases/pp/Gibbons.htm}{LSU Law Center}).

\textbf{Rule Announced:} Under the Commerce Clause (U.S. Const. art. I, \S~8, cl. 3), ``commerce'' encompasses not merely buying and selling commodities but every species of commercial intercourse, including navigation by steam-propelled vessels carrying passengers for hire (\href{https://biotech.law.lsu.edu/cases/pp/Gibbons.htm}{LSU Law Center}; \href{https://www.law.cornell.edu/wex/enumerated_powers}{Cornell LII}). Commerce ``among the several States'' reaches commercial intercourse concerning more than one state and does not stop at state boundaries, though it does not extend to commerce that is ``completely internal'' to a single state (\href{https://biotech.law.lsu.edu/cases/pp/Gibbons.htm}{LSU Law Center}). Congress's power to regulate such commerce is plenary, complete in itself, and subject only to constitutional limits. Under the Supremacy Clause (U.S. Const. art. VI, cl. 2), a federal coasting license issued under the Coasting Act of 1793 affirmatively authorizes interstate navigation and preempts any conflicting state monopoly statute (\href{https://biotech.law.lsu.edu/cases/pp/Gibbons.htm}{LSU Law Center}; \href{https://www.courtlistener.com/opinion/85412/gibbons-v-ogden/}{CourtListener}).

\textbf{Holding vs. Dicta:} The binding \emph{holding} rests on statutory preemption grounded in Congress's affirmative Commerce Clause power: navigation is commerce; Congress validly enacted the Coasting Act of 1793; a federal coasting license confers an affirmative right to navigate interstate waters regardless of propulsion method; and New York's monopoly directly collided with that federal statute and therefore had to yield under the Supremacy Clause (\href{https://biotech.law.lsu.edu/cases/pp/Gibbons.htm}{LSU Law Center}). Several celebrated passages are \emph{dicta}: (1) Marshall's sympathetic discussion of Webster's ``exclusivity'' (dormant Commerce Clause) theory---stating ``[t]here is great force in this argument, and the Court is not satisfied that it has been refuted''---which Marshall expressly declined to decide because Congress had already acted (\href{https://biotech.law.lsu.edu/cases/pp/Gibbons.htm}{LSU Law Center}; \href{https://scholarship.law.slu.edu/cgi/viewcontent.cgi?article=1753&context=lj}{Norman R. Williams, St. Louis U. L.J.}); (2) Marshall's description of ``that immense mass of legislation'' (inspection, quarantine, and health laws, internal commerce, turnpikes, and ferries) remaining within state police power; and (3) the Patent and Copyright Clause (art. I, \S~8, cl. 8) challenge, which the Court expressly deemed unnecessary to resolve (\href{https://biotech.law.lsu.edu/cases/pp/Gibbons.htm}{LSU Law Center}).

\section{Analysis of Opinions}
\subsection{Majority Opinion (Chief Justice Marshall)}
Chief Justice Marshall began by rejecting Ogden's plea for a ``strict construction'' of federal enumerated powers, reasoning that the Constitution contains no text mandating a narrow construction that would ``cripple the government'' and that constitutional terms should receive their natural, obvious meaning in light of the instrument's objects (\href{https://biotech.law.lsu.edu/cases/pp/Gibbons.htm}{LSU Law Center}; \href{https://teachingamericanhistory.org/document/gibbons-v-ogden/}{Teaching American History}).

Marshall then construed each phrase of the Commerce Clause:
\begin{itemize}
\item \textbf{``Commerce'':} Rejecting Ogden's attempt to confine ``commerce'' to mere ``traffic'' (buying and selling commodities), Marshall defined commerce as ``commercial intercourse'' in all its branches, expressly encompassing navigation. He reinforced this reading with constitutional text (the Port Preference Clause of art. I, \S~9, cl. 6, which presupposes federal power over vessels) and historical practice (including federal embargoes and navigation acts) (\href{https://biotech.law.lsu.edu/cases/pp/Gibbons.htm}{LSU Law Center}).
\item \textbf{``Among the several States'':} Defining ``among'' as ``intermingled with,'' Marshall concluded that federal authority cannot stop at state border lines but extends into a state's interior whenever commerce concerns more than one state. Conversely, the enumeration presupposes something not enumerated: the ``exclusively internal commerce of a State'' that does not affect other states (\href{https://biotech.law.lsu.edu/cases/pp/Gibbons.htm}{LSU Law Center}).
\item \textbf{``To regulate'':} Marshall defined regulation as prescribing the rule by which commerce is governed---a plenary power, complete in itself, whose sole restraints (absent explicit constitutional prohibitions) are political checks through representative elections (\href{https://teachingamericanhistory.org/document/gibbons-v-ogden/}{Teaching American History}).
\end{itemize}

Turning to state power, Marshall distinguished concurrent state taxation from commercial regulation and explained that state inspection, quarantine, and health laws flow from state police powers rather than a concurrent power over interstate commerce (\href{https://biotech.law.lsu.edu/cases/pp/Gibbons.htm}{LSU Law Center}). Rather than deciding whether the Commerce Clause alone invalidated New York's monopoly in the absence of federal legislation, Marshall pivoted to the Supremacy Clause: even if New York acted under an acknowledged state police power, its law had to yield to a conflicting federal statute. Construing \S\S~1 and 4 of the Coasting Act of 1793 (together with an 1812 steamboat enrollment statute), Marshall held that a federal coasting license grants affirmative permission to carry on the coasting trade---not merely American national character---whether the vessel carries cargo or passengers and whether propelled by wind or steam (\href{https://biotech.law.lsu.edu/cases/pp/Gibbons.htm}{LSU Law Center}). Because New York's monopoly directly collided with Gibbons's federal licenses, the state statutes were void to that extent.

\subsection{Concurring Opinion (Justice Johnson)}
Justice William Johnson concurred in the judgment but wrote separately to ground the reversal directly on the Constitution rather than on statutory preemption under the Coasting Act (\href{https://teachingamericanhistory.org/document/gibbons-v-ogden/}{Teaching American History}; \href{https://www.oyez.org/cases/1789-1850/22us1}{Oyez}). Reasoning from the law of nations and the historical purpose of the Constitution to end interstate trade wars, Johnson argued that the power to regulate interstate and foreign commerce was transferred in its entirety from the states to Congress. Because the power to prescribe limits on commerce necessarily includes the power to decide what shall remain unrestrained, the Commerce Clause is strictly \emph{exclusive}---``it can reside but in one potentate''---leaving states with zero authority to regulate interstate navigation even if Congress had never enacted the Coasting Act (\href{https://teachingamericanhistory.org/document/gibbons-v-ogden/}{Teaching American History}; \href{https://scholarship.law.slu.edu/cgi/viewcontent.cgi?article=1753&context=lj}{Williams, St. Louis U. L.J.}).

\subsection{Dissenting Opinions}
No Justice dissented. Justice Smith Thompson did not participate in the consideration or decision of the case (\href{https://www.oyez.org/cases/1789-1850/22us1}{Oyez}).

\section{Examples: Future Applications}
\begin{enumerate}
\item \textbf{Same-Side (Federal Power Upheld / State Monopoly Invalidated):} Congress enacts a federal statute licensing interstate commercial passenger ferries and freight barges operating on the Mississippi River between Illinois and Missouri. Illinois grants a private corporation an exclusive twenty-year franchise to operate electric-powered vessels in Illinois waters and sues to enjoin a federally licensed Missouri operator. Under \textit{Gibbons}, interstate water transportation of passengers or cargo is ``commerce among the several States,'' the federal license confers affirmative authority regardless of propulsion technology, and Illinois's conflicting monopoly is preempted under the Supremacy Clause.
\item \textbf{Same-Side (Federal Power Upheld / State Restriction Invalidated):} Congress passes a comprehensive statute licensing interstate commercial trucking carriers to operate on the federal highway network. A state bars federally licensed autonomous freight trucks from entering its borders in order to protect in-state trucking firms. Just as New York could not exclude steam-propelled vessels licensed under the Coasting Act, the state's restriction directly collides with a valid federal regulation of interstate commercial intercourse and must yield.
\item \textbf{Opposite-Side (State Regulation Upheld / Beyond Federal Preemption):} A state enacts a non-discriminatory public health quarantine and agricultural inspection law requiring livestock raised within the state to pass disease inspection before entering local markets, where Congress has enacted no conflicting federal statute. Under Marshall's distinction in \textit{Gibbons}, inspection and health laws acting on local goods prior to interstate transit fall within the state's reserved ``immense mass of legislation'' (police power) and do not intrude upon federal commercial authority.
\item \textbf{Opposite-Side (Purely Internal Commerce / Outside \textit{Gibbons}'s Direct Scope):} A state regulates licensing and safety standards for small rowboats rented exclusively for recreational use on a landlocked lake situated wholly within a single county, with no connection to out-of-state waterways or interstate travelers. Under Marshall's boundary in \textit{Gibbons}, commerce ``completely internal'' to one state that ``does not extend to or affect other States'' remains reserved to the state itself.
\item \textbf{Fence-Sitter (Ambiguous / Contested Boundary):} A state grants an exclusive franchise to operate a high-speed commuter ferry running strictly between two ports within the same state (for example, Albany and New York City along the Hudson River), but the ferry carries many out-of-state business travelers and operates on navigable waters heavily used by interstate vessels, and Congress has enacted no statute specifically governing intrastate routes. Under \textit{Gibbons}, the outcome is uncertain: Marshall reserved ``completely internal'' commerce to the states and rested his holding on actual statutory collision rather than pure dormancy, whereas Johnson's exclusivity theory---and Marshall's recognition that Congress's power extends into a state's interior when commerce ``concerns more States than one''---would push toward invalidation if the franchise burdens interstate navigation.
\end{enumerate}

\section{Critique}
\textbf{Originalist and Textualist Perspectives:} Originalist scholars sharply debate how faithfully Marshall's definitions tracked 1787 public meaning. Randy E. Barnett has argued from extensive founding-era evidence that ``commerce'' originally meant trade and transportation of goods (including navigation), supporting Marshall's result in \textit{Gibbons} while rejecting later twentieth-century expansions of ``commerce'' to manufacturing or agriculture (\href{https://scholarship.law.georgetown.edu/context/facpub/article/1504/viewcontent/barnett_commerce_clause.pdf}{Barnett, \textit{The Original Meaning of the Commerce Clause}}; \href{https://constitutioncenter.org/the-constitution/articles/article-i/clauses/752}{National Constitution Center}). By contrast, Jack M. Balkin's ``living originalism'' reads Marshall's phrase ``intercourse'' as embracing broader social and economic ``interaction'' across state lines, a reading Barnett and Heritage Foundation commentators contest by noting that founding-era ``commercial intercourse'' still denoted cross-border movement and exchange rather than all human interaction (\href{https://constitution.heritage.org/essays/a1-s8-c3-b/}{Heritage Guide to the Constitution}; \href{https://constitutioncenter.org/the-constitution/articles/article-i/clauses/752}{National Constitution Center}). Textualist critics also question Marshall's expansive gloss on ``among'' (``intermingled with'') and his dicta endorsing an exclusive, unenumerated Dormant Commerce Clause, noting that the constitutional text grants an affirmative power to Congress without expressly stripping states of concurrent authority (\href{https://constitution.heritage.org/essays/a1-s8-c3-b/}{Heritage Guide}; \href{https://www.uclalawreview.org/wp-content/uploads/2019/09/11_53UCLALRev1532005-2006.pdf}{Williams, UCLA L. Rev.}).

\textbf{Progressive and Doctrinal Perspectives:} Doctrinal and progressive scholars celebrate \textit{Gibbons} as the architectural cornerstone of the American common market and modern national regulatory authority, enabling Congress to address nationwide economic, labor, and civil rights problems that states cannot solve individually (\href{https://www.archives.gov/milestone-documents/gibbons-v-ogden}{National Archives}; \href{https://www.law.cornell.edu/wex/enumerated_powers}{Cornell LII}). Yet doctrinal historians such as Norman R. Williams have critiqued Marshall's strategic ``hedge'': after spending pages dismantling concurrent state commerce power, Marshall pivoted to a strained reading of the Coasting Act of 1793 to avoid a pure Dormant Commerce Clause holding, leaving lower courts and subsequent Supreme Courts mired in decades of doctrinal confusion over state commercial regulation until \textit{Cooley v. Board of Wardens} (1851) (\href{https://scholarship.law.slu.edu/cgi/viewcontent.cgi?article=1753&context=lj}{Williams, St. Louis U. L.J.}; \href{https://www.nyulawreview.org/wp-content/uploads/2018/08/NYULawReview-79-4-Williams.pdf}{Williams, N.Y.U. L. Rev.}; \href{https://history.nycourts.gov/wp-content/uploads/2026/05/25_Judicial_Notice-compressed.pdf}{Historical Society of the New York Courts}). Moreover, Marshall's dual commitments---declaring Congress's power ``complete in itself'' while simultaneously insisting that ``completely internal'' commerce belongs to the states---created an enduring tension that resurfaced in \textit{Wickard v. Filburn} (1942), \textit{United States v. Lopez} (1995), and \textit{Gonzales v. Raich} (2005) (\href{https://www.law.cornell.edu/wex/enumerated_powers}{Cornell LII}; \href{https://www.everycrsreport.com/reports/RL32844.html}{CRS Report RL32844}).

\textbf{Underlying Competing Values:} \textit{Gibbons} sits at the fault line between \emph{national economic integration} (preventing protectionist state monopolies and retaliatory trade wars like those between New York, New Jersey, and Connecticut) and \emph{federalist decentralization} (preserving state sovereignty, local police powers, and democratic experimentation within enumerated federal limits).

\section{Key Quotations}
\begin{quote}
``Commerce, undoubtedly, is traffic, but it is something more: it is intercourse. It describes the commercial intercourse between nations, and parts of nations, in all its branches, and is regulated by prescribing rules for carrying on that intercourse.''
\end{quote}
\noindent --- Chief Justice John Marshall, 22 U.S. at 189--90 (\href{https://biotech.law.lsu.edu/cases/pp/Gibbons.htm}{LSU Law Center}; \href{https://www.courtlistener.com/opinion/85412/gibbons-v-ogden/}{CourtListener}).

\begin{quote}
``The word `among' means intermingled with. A thing which is among others is intermingled with them. Commerce among the States cannot stop at the external boundary line of each State, but may be introduced into the interior.''
\end{quote}
\noindent --- Chief Justice John Marshall, 22 U.S. at 194 (\href{https://teachingamericanhistory.org/document/gibbons-v-ogden/}{Teaching American History}; \href{https://www.courtlistener.com/opinion/85412/gibbons-v-ogden/}{CourtListener}).

\begin{quote}
``The enumeration presupposes something not enumerated; and that something, if we regard the language or the subject of the sentence, must be the exclusively internal commerce of a State.''
\end{quote}
\noindent --- Chief Justice John Marshall, 22 U.S. at 195 (\href{https://biotech.law.lsu.edu/cases/pp/Gibbons.htm}{LSU Law Center}; \href{https://www.courtlistener.com/opinion/85412/gibbons-v-ogden/}{CourtListener}).

\begin{quote}
``It is the power to regulate; that is, to prescribe the rule by which commerce is to be governed. This power, like all others vested in Congress, is complete in itself, may be exercised to its utmost extent, and acknowledges no limitations, other than are prescribed in the Constitution.''
\end{quote}
\noindent --- Chief Justice John Marshall, 22 U.S. at 196 (\href{https://teachingamericanhistory.org/document/gibbons-v-ogden/}{Teaching American History}; \href{https://www.courtlistener.com/opinion/85412/gibbons-v-ogden/}{CourtListener}).

\begin{quote}
``And since the power to prescribe the limits to its freedom necessarily implies the power to determine what shall remain unrestrained, it follows that the power must be exclusive; it can reside but in one potentate, and hence the grant of this power carries with it the whole subject, leaving nothing for the State to act upon.''
\end{quote}
\noindent --- Justice William Johnson, concurring in the judgment, 22 U.S. at 227 (\href{https://teachingamericanhistory.org/document/gibbons-v-ogden/}{Teaching American History}; \href{https://www.courtlistener.com/opinion/85412/gibbons-v-ogden/}{CourtListener}).

\end{document}
